% Reporte de la Actividad 5
% Exclusión mutua con Test-and-Set en Java y C
\documentclass[12pt,oneside]{template/liverpoolthesis}
\input{template/preamble}

% ---------- Datos de esta actividad ----------
\title{Actividad 5: Exclusión mutua con Test-and-Set}
\author{%
  Andrea Hernandez Huerta\\
  ID: 175523\\[1em]
  Heriberto Espino Montelongo\\
  ID: 175199\\[1em]
  \text{Profesor:} Dr. Victor Giovanni Morales Murillo}
\faculty{Universidad de las Américas Puebla (UDLAP)}
\school{Licenciatura en Ciencia de Datos}
\degree{Curso: Programación Concurrente}
\date{Fecha de entrega: 17 de septiembre de 2026}
\subject{Exclusión mutua con Test-and-Set}

\addto\captionsspanish{\renewcommand{\chaptername}{Ejercicio}}
\providecommand{\figureautorefname}{Figura}
\addto\captionsspanish{\renewcommand{\figureautorefname}{Figura}}

\begin{document}

\frontmatter
\maketitle
\tableofcontents
\listoffigures
\lstlistoflistings

\mainmatter

\chapter{Exclusión mutua con Test-and-Set}

\textbf{Desarrollen un programa en Java que cree un cerrojo utilizando
\texttt{AtomicBoolean}, inicialmente en \texttt{false}; cree tres hilos que
compartan el mismo cerrojo y un contador; implemente una operación para
solicitar acceso a la sección crítica utilizando \texttt{getAndSet(true)};
mantenga al hilo en espera activa mientras el cerrojo esté ocupado; incremente
el contador y muestre el identificador y el valor del contador dentro de la
sección crítica; libere el cerrojo cambiándolo a \texttt{false}; y repita el
proceso cinco veces por hilo.}

\emph{Entradas.} Tres hilos identificados como $P0$, $P1$ y $P2$, un
\texttt{AtomicBoolean lockFlag} compartido e inicialmente en \texttt{false},
un contador compartido inicializado en cero y cinco iteraciones por hilo. La
prueba de inicio permite que los tres hilos comiencen a solicitar el cerrojo
de manera concurrente.

\emph{Procesamiento.} La sección de entrada ejecuta
\texttt{requestCS}. En Java, la expresión \texttt{lockFlag.getAndSet(true)}
obtiene el valor anterior y cambia el cerrojo a \texttt{true} como una sola
operación atómica. Si el valor anterior es \texttt{true}, el hilo permanece en
el ciclo \texttt{while} y realiza espera activa. Si el valor anterior es
\texttt{false}, el hilo entra a la sección crítica, incrementa el contador,
muestra su identificador y libera el cerrojo con \texttt{getAndSet(false)}.
La implementación en C conserva la misma estructura mediante
\texttt{atomic\_exchange} sobre un \texttt{atomic\_bool}.

\emph{Salidas.} Cada entrada muestra el hilo, la iteración y el valor del
contador. Como existen tres hilos y cada uno realiza cinco entradas, el valor
esperado es $3\times5=15$. El valor obtenido en las pruebas fue $15$ tanto en
Java como en C.

\emph{Estructuras de control utilizadas.} Se emplean tres hilos, un ciclo
\texttt{for} para repetir las cinco entradas, un ciclo \texttt{while} para la
espera activa, métodos de solicitud y liberación del cerrojo y una sección
crítica entre las secciones de entrada y salida. También se utilizan
operaciones condicionales implícitas en el valor anterior que devuelve
\texttt{getAndSet} o \texttt{atomic\_exchange}.

\emph{Conceptos.} Test-and-Set es una instrucción capaz de leer una variable y
modificarla como una única operación atómica. Al convertir la comprobación y
la modificación del cerrojo en una sola operación, solo un hilo puede
obtener \texttt{false} cuando el cerrojo está libre. Los demás hilos obtienen
\texttt{true} y permanecen en espera activa hasta que otro hilo libera el
cerrojo. La espera activa consume tiempo de procesador, que es la limitación
principal indicada en el módulo.

\emph{Secciones del programa.}
\begin{itemize}
  \item \textbf{Sección de entrada:} el método \texttt{requestCS} ejecuta
  repetidamente Test-and-Set.
  \item \textbf{Sección crítica:} el hilo incrementa el contador y muestra
  el resultado.
  \item \textbf{Sección de salida:} el método \texttt{releaseCS} cambia el
  cerrojo a \texttt{false}.
  \item \textbf{Sección restante:} la repetición de la prueba y la espera de
  los hilos al finalizar.
\end{itemize}

\emph{Prueba de escritorio.} La siguiente secuencia representa el caso en
que $P0$ ya obtuvo el cerrojo y $P1$ intenta entrar simultáneamente. Los
valores se expresan como booleanos, de acuerdo con el enunciado.

\begin{center}
\scriptsize
\resizebox{\textwidth}{!}{%
\begin{tabular}{c|c|c|c|c|p{4.0cm}}
Paso & Proceso & \texttt{lockFlag} antes & Operación & \texttt{lockFlag} después & Resultado \\
\hline
1 & $P0$ & \texttt{false} & \texttt{TestAndSet(true)} devuelve \texttt{false} & \texttt{true} & $P0$ entra a la sección crítica. \\
2 & $P1$ & \texttt{true} & \texttt{TestAndSet(true)} devuelve \texttt{true} & \texttt{true} & $P1$ no entra y comienza la espera activa. \\
3 & $P1$ & \texttt{true} & \texttt{TestAndSet(true)} devuelve \texttt{true} & \texttt{true} & $P1$ continúa esperando. \\
4 & $P0$ & \texttt{true} & \texttt{TestAndSet(false)} & \texttt{false} & $P0$ termina y libera el cerrojo. \\
5 & $P1$ & \texttt{false} & \texttt{TestAndSet(true)} devuelve \texttt{false} & \texttt{true} & $P1$ obtiene el cerrojo y entra. \\
6 & $P1$ & \texttt{true} & \texttt{TestAndSet(false)} & \texttt{false} & $P1$ termina y libera el cerrojo. \\
\end{tabular}
}%
\end{center}

\emph{Evaluación.} El algoritmo garantiza exclusión mutua porque la operación
que obtiene el valor anterior y cambia el cerrojo es una operación atómica
desde el punto de vista de los hilos que compiten. Cuando el cerrojo está libre, solo
un hilo puede recibir \texttt{false}; los demás reciben \texttt{true} y no
avanzan a la sección crítica. El contador final coincide con el esperado: 15.

\emph{Espera activa.} La espera activa aparece en el ciclo
\texttt{while} que contiene \texttt{getAndSet(true)} de Java y en el ciclo
equivalente \texttt{while} que contiene \texttt{testAndSet(true)} de C. En cada
repetición el hilo vuelve a comprobar el cerrojo y cede brevemente el
procesador, pero permanece listo para ejecutar. Esta solución evita la entrada
simultánea, aunque consume CPU mientras existen hilos esperando.

\section{Código fuente en Java}

El siguiente listado conserva la forma del pseudocódigo del módulo: la
solicitud usa \texttt{getAndSet(true)}, la liberación cambia el cerrojo a
\texttt{false} y el contador se modifica dentro de la sección crítica.

\lstinputlisting[
  language=Java,
  caption={Implementación de Test-and-Set en Java},
  label={lst:java-test-and-set}
]{../code/java/Ejercicio01TestAndSet.java}

\section{Código equivalente en C}

En C, \texttt{atomic\_exchange} realiza la misma operación de lectura y
modificación sobre \texttt{atomic\_bool lockFlag}. La función
\texttt{testAndSet} conserva el nombre conceptual del pseudocódigo y devuelve
el valor anterior del cerrojo.

\lstinputlisting[
  language=C,
  caption={Implementación equivalente de Test-and-Set en C},
  label={lst:c-test-and-set}
]{../code/c/Ejercicio01TestAndSet.c}

\section{Registro de ejecución}

La salida depende del orden en que el planificador concede el procesador a
cada hilo. Lo importante es que cada entrada ocurre con un valor distinto y
que el contador termina en 15.

\begin{figure}[H]
\centering
\begin{minipage}{0.94\textwidth}
\begin{lstlisting}[language={},numbers=none]
prueba de exclusion mutua con test-and-set
Hilo 2 entra a la seccion critica en la iteracion 1, contador = 1
Hilo 1 entra a la seccion critica en la iteracion 1, contador = 2
Hilo 0 entra a la seccion critica en la iteracion 1, contador = 3
Hilo 1 entra a la seccion critica en la iteracion 2, contador = 4
Hilo 0 entra a la seccion critica en la iteracion 2, contador = 5
Hilo 2 entra a la seccion critica en la iteracion 2, contador = 6
...
contador esperado: 15
contador obtenido: 15
resultado: getAndSet garantiza exclusion mutua con un cerrojo booleano
\end{lstlisting}
\end{minipage}
\caption{Registro de ejecución de la implementación en Java.}
\label{fig:ejecucion-test-and-set}
\end{figure}

\chapter{Conclusiones}

La implementación demuestra la función de Test-and-Set dentro de la
exclusión mutua. El cerrojo booleano inicia libre, cada hilo solicita acceso
con una única operación que obtiene el valor anterior y escribe
\texttt{true}, y el hilo que logra entrar libera el cerrojo al salir. Por
ello, tres hilos que realizan cinco entradas producen exactamente quince
incrementos del contador.

La prueba de escritorio muestra la diferencia con una solución que primero
comprueba y después modifica el cerrojo. Con Test-and-Set esas dos acciones se
realizan como una operación atómica, así que dos hilos no pueden obtener
\texttt{false} para la misma entrada. Su costo es la espera activa, que
consume CPU mientras el cerrojo permanece ocupado.

\cleardoublepage
\phantomsection
\addcontentsline{toc}{chapter}{Referencias}
\nocite{*}
\bibliography{template/references}

\end{document}
